随着高校毕业论文和学术论文写作需求的不断增加，论文排版的规范性和效率逐渐成为学生和科研人员关注的重要问题。传统文字处理工具虽然操作简单，但在复杂公式、图表编号、参考文献管理和模板适配等方面存在一定不足；LaTeX 具有较强的学术排版能力，但其语法规则、编译环境和错误日志对初学者而言具有较高门槛。针对上述问题，本文设计并实现了一套基于 JavaWeb 和 LaTeX 的学术论文智能排版系统，旨在为用户提供在线化、模板化和智能化的论文写作与排版支持。

系统采用前后端分离架构，前端基于 Vue 3 和 Vite 构建用户交互界面，后端基于 Spring Boot、MyBatis 和 MySQL 实现业务逻辑处理与数据持久化。系统围绕论文写作流程，主要实现了用户注册登录、邮箱验证码、个人资料管理、论文项目管理、LaTeX 在线编辑、模板管理、论文编译、PDF 预览、文件导出、AI 对话辅助和编译错误分析等功能。用户可以在浏览器中创建论文项目，选择或导入 LaTeX 模板，在线编辑论文源码，并通过多种编译引擎生成 PDF 文档。

在系统实现过程中，针对 LaTeX 模板依赖文件较多、编译错误不易理解等问题，系统支持 zip 模板包导入，并在编译过程中尽量保留模板目录结构，减少因 .cls、.sty、.bib 等依赖文件缺失导致的编译失败。同时，系统引入 AI 辅助功能，在用户遇到 LaTeX 语法问题或编译失败时，能够根据错误日志给出原因分析和修改建议，从而降低 LaTeX 使用难度。

测试结果表明，该系统能够较好地完成论文项目管理、模板载入、在线编辑、论文编译、PDF 预览和 AI 错误分析等任务，基本满足学术论文在线排版的使用需求，具有一定的实用价值和应用前景。